\chapter*{Abstract}
\justify
India carries one of the largest diabetes burdens in the world, and when an adult with type 2 diabetes is not controlled on metformin, the clinician must choose a second-line drug with little patient-specific evidence. Existing digital tools either restate guideline rules, predict risk from cross-sectional data such as the Pima dataset, or answer free-text questions through general-purpose language models that can produce uncited or unsafe content. Predictive models answer ``who is at risk'' but not ``what would happen to this patient under each drug'', which is a causal question distorted by confounding by indication. This work presents DiaCausal, a research prototype that first removes unsafe options with deterministic, cited safety rules, then estimates the six-month HbA1c change under a sodium--glucose co-transporter-2 inhibitor, a dipeptidyl peptidase-4 inhibitor and a sulfonylurea using a cross-fitted multinomial propensity model, augmented inverse probability weighting and a doubly robust learner, each with a 95\% interval, and abstains when overlap is insufficient. Because no open Indian dataset records treatment and outcome, the engine is developed and benchmarked on an India-calibrated synthetic cohort with known true effects. Over 20 synthetic cohorts of 5{,}000 patients, the naive comparison overstated the SGLT2i-versus-DPP-4i effect by about 75\%, whereas AIPW reduced the bias to at most 0.005 HbA1c points with 90--95\% interval coverage, and the DR-learner's patient-level intervals covered the truth for 92.7--94.7\% of test patients. The engine runs in a live, installable phone website with admin-approved accounts and an early retrieval layer that returns cited passages from licence-cleared sources. This mid-semester report presents the proposed system, requirements, design, timeline, the implementation completed so far, initial results and the evaluation matrix. Intended use: \intendeduse
